\section{Noetherian Rings}

\subsection{Exercise - Complete}\label{7.1}

\begin{proof}
Let $\Sigma$ be the set of ideals in $A$ which are not finitely generated. Since $A$ is not Noetherian, $\Sigma$ is non-empty. Let $\mathcal{T}$ be a chain
in $\Sigma$. Define $\mathfrak{a} = \bigcup_{\mathfrak{b} \in \mathcal{T}}\mathfrak{b}$. If $x_1,x_2 \in \mathfrak{a}$, then $x_1 \in \mathfrak{b}_1$ and $x_2 \in \mathfrak{b}_2$
for ideals $\mathfrak{b}_1,\mathfrak{b}_2 \in \mathcal{T}$. Since $\mathcal{T}$ is a chain, these ideals are comparable, so relabelling if necessary, we may assume that
$\mathfrak{b}_1 \subset \mathfrak{b}_2$. Then, $x_1,x_2 \in \mathfrak{b}_2$, so $x_1 + x_2 \in \mathfrak{b}_2 \subset \mathfrak{a}$. Thus, $\mathfrak{a}$ is additively closed.
If $a \in A$ and $x \in \mathfrak{a}$, there is some $\mathfrak{b} \in \mathcal{T}$ such that $x \in \mathfrak{b}$. Then, $ax \in \mathfrak{b} \subset \mathfrak{a}$, so
$\mathfrak{a}$ is an ideal of $A$.

Suppose to the contrary that $\mathfrak{a}$ is finitely generated. Then, there are finitely many generating elements $x_1,\ldots,x_n$. Each $x_i$ is in some $\mathfrak{b}_i \in \mathcal{T}$.
Since $\mathcal{T}$ is a chain, there is some $j \in \{1,\ldots,n\}$ such that $\mathfrak{b}_i \subset \mathfrak{b}_j$ for all $i$. Then, $x_i \in \mathfrak{b}_i \subset \mathfrak{b}_j$
for all $i$. Since $\mathfrak{b}_j \subset \mathfrak{a}$ and $\mathfrak{a} = \sum_i A x_i \subset \sum_i \mathfrak{b}_j = \mathfrak{b}_j$, we have $\mathfrak{a} = \mathfrak{b}_j$. But
since $\mathfrak{b}_j \in \Sigma$, $\mathfrak{a} = \mathfrak{b}_j$ is not finitely generated, a contradiction. Thus, $\mathfrak{a} \in \Sigma$, so all chains in $\Sigma$ have maximal
elements in $\Sigma$. By Zorn's lemma, $\Sigma$ has maximal elements.

Now, suppose $\mathfrak{a}$ is a maximal element of $\Sigma$. Suppose there exist $x \not\in \mathfrak{a}$ and $y \not\in \mathfrak{a}$ such that $xy \in \mathfrak{a}$. Thus, the ideal
$\mathfrak{a} + (x)$ is an ideal that properly contains $\mathfrak{a}$. Since $\mathfrak{a}$ is maximal in $\Sigma$, $\mathfrak{a} + (x)$ is finitely generated. In particular, there are
elements $x_1,\ldots,x_n$ that generate $\mathfrak{a} + (x)$ as an $A$-module. Then, for each $x_i$, there is some $a_i \in A$ such that $x_i - a_i x \in \mathfrak{a}$. Let $\mathfrak{a}_0$
be the ideal generated by $\{x_i - a_i x\}_i$, so clearly, $\mathfrak{a}_0 \subset \mathfrak{a}$ and $\mathfrak{a}_0$ is finitely generated. Further, $x_i \in \mathfrak{a}_0 + (x)$ for all $i$
since $x_i - a_i x \in \mathfrak{a}_0$ and $a_i x \in (x)$. Therefore, $\mathfrak{a}_0 + (x) = \mathfrak{a} + (x)$.

By the definition of $(\mathfrak{a}:x)$, $x(\mathfrak{a}:x) \subset \mathfrak{a}$. Further, $\mathfrak{a} \subset (\mathfrak{a}:x)$; this inclusion is proper since $y \in (\mathfrak{a}:x)$ but
$y \not\in \mathfrak{a}$. Since $\mathfrak{a}$ is maximal in $\Sigma$, $(\mathfrak{a}:x)$ is finitely generated; hence, $x(\mathfrak{a}:x)$ is finitely generated. Therefore, 
$\mathfrak{a}_0 + x(\mathfrak{a}:x) \subset \mathfrak{a}$. Suppose $a \in \mathfrak{a}$. Then, $a = a' + bx$, where $a' \in \mathfrak{a}_0$ and $b \in A$ 
since $\mathfrak{a}_0 + (x) = \mathfrak{a} + (x)$. Thus, $bx = a - a' \subset \mathfrak{a}$, so $b \in (\mathfrak{a}:x)$. Hence, $bx \in x(\mathfrak{a}:x)$ so that 
$a = a' + bx \in \mathfrak{a}_0 + x(\mathfrak{a}:x)$. Therefore, we have equality: $\mathfrak{a}_0 + x(\mathfrak{a}:x) = \mathfrak{a}$. 

Since $\mathfrak{a}_0$ and $x(\mathfrak{a}:x)$ are both finitely generated, their sum, $\mathfrak{a}$, is also finitely generated. This is a clear contradiction,
so we cannot have $x \not\in \mathfrak{a}$ and $y \not\in \mathfrak{a}$ such that $xy \in \mathfrak{a}$; hence, $\mathfrak{a}$ is prime. Thus, maximal elements of $\Sigma$ is prime.
In particular, if $A$ is a ring in which all prime ideals are finitely generated, then $\Sigma$ must be empty, so every ideal of $A$ is finitely generated; in other words, $A$
is Noetherian.
\end{proof}

\

\subsection{Exercise - Complete}\label{7.2}

\begin{proof}
If $\sum_i a_i x^i \in A[[x]]$ is nilpotent, then by Exercise \ref{1.5}, each $a_i \in A$ is nilpotent. Conversely, suppose $a_i \in A$ is nilpotent for
all $i$. Since $A$ is Noetherian, we have by Corollary 7.15 that the nilradical is nilpotent. Thus, $\mathfrak{N}^m = 0$ for some $m > 0$. Expanding 
$(\sum_i a_i x^i)^m$, we obtain a sum of elements of the form $a_{i_1} \cdots a_{i_m} x^{i_1 + \cdots + i_m}$. Since $a_{i_k} \in \mathfrak{N}$
for all $k$, $a_{i_1} \cdots a_{i_m} \in \mathfrak{N}^m = 0$ so that each term in the sum vanishes. Thus, $(\sum_i a_i x^i)^m = 0$, as desired.
\end{proof}

\subsection{Exercise}\label{7.3}

\subsection{Exercise}\label{7.4}

\subsection{Exercise}\label{7.5}

\subsection{Exercise}\label{7.6}

\subsection{Exercise}\label{7.7}

\subsection{Exercise}\label{7.8}

\subsection{Exercise}\label{7.9}

\subsection{Exercise}\label{7.10}

\subsection{Exercise}\label{7.11}

\subsection{Exercise}\label{7.12}

\subsection{Exercise}\label{7.13}

\subsection{Exercise}\label{7.14}

\subsection{Exercise}\label{7.15}

\subsection{Exercise}\label{7.16}

\subsection{Exercise}\label{7.17}

\subsection{Exercise}\label{7.18}

\subsection{Exercise}\label{7.19}

\subsection{Exercise}\label{7.20}

\subsection{Exercise}\label{7.21}

\subsection{Exercise}\label{7.22}

\subsection{Exercise}\label{7.23}

\subsection{Exercise}\label{7.24}

\subsection{Exercise}\label{7.25}

\subsection{Exercise}\label{7.26}

\subsection{Exercise}\label{7.27}


